%=========== ΛΥΣΕΙΣ - ΔΙΑΓΩΝΙΣΜΑ ΠΡΟΣΟΜΟΙΩΣΗΣ 15 ===========

%=========== ΛΥΣΗ - 15ο ΘΕΜΑ Α (Σ/Λ μόνο) ===========
\begin{thema}{Α}
\begin{erwthma}
\setcounter{enumi}{3}
\item \ \\
\begin{alist}
\item \textbf{Σ.} 
Είναι γνωστό και θεμελιώδες συμπέρασμα. Η γνήσια μονοτονία αποτρέπει ποτέ δύο διαφορετικά $x$ να επιστρέψουν το ίδιο $y$, εξασφαλίζοντας τέλεια τη συνθήκη $1-1$.
\item \textbf{Λ.}
Το $\dint_\alpha^\beta f(x)\,\d x$ είναι μια απόλυτα \textbf{σταθερή αριθμητική τιμή} (ένας συγκεκριμένος πραγματικός αριθμός, εκφράζοντας γεωμετρικό εμβαδόν). Επομένως, το διαφορικό (παράγωγος) οποιουδήποτε σταθερού αριθμού ισούται με μηδέν, δηλαδή $\left( \dint_\alpha^\beta f(x)\,\d x \right)' = 0$ και όχι $f(x)$. (Η προσοχή βρίσκεται στη διαφορά αόριστου συνόρου ολοκλήρωσης από ορισμένο).
\item \textbf{Σ.}
Πράγματι, εάν δεν υπάρχουν αρνητικές περιοχές να ακυρώσουν θετικές (αφού $f(x) \geq 0$) και το εμβαδόν είναι μηδέν, η συνάρτηση (αν είναι συνεχής) πρέπει να ταυτίζεται αυστηρά με τον άξονα $x'x$ (δηλαδή $f(x)=0$) σε όλο το διάστημα.
\item \textbf{Σ.}
Η σύνθεση διατηρεί τη μονοσήμαντη αντιστοίχιση.
\item \textbf{Σ.}
Πρόκειται για το άμεσο απόρροια του Θεωρήματος Μέγιστης-Ελάχιστης Τιμής για συνεχείς συναρτήσεις ορισμένες σε περίκλειστα διαστήματα.
\end{alist}
\end{erwthma}
\end{thema}

%=========== ΛΥΣΗ - 15ο ΘΕΜΑ Β ===========
\begin{thema}{Β}
Δίνεται $f(x) = \frac{x^2-4}{x+2}$ με $x \neq -2$.
\begin{erwthma}

\item \textbf{Απλοποίηση και Όρια:}
Παραγοντοποιούμε τη διαφορά τετραγώνων στον αριθμητή:
\[ f(x) = \frac{(x-2)(x+2)}{x+2} \]
Δεδομένου ότι το πεδίο ορισμού εξαιρεί το $x=-2$, η παρένθεση $(x+2)$ δεν είναι μηδέν. Μπορούμε να ολοκληρώσουμε την απλοποίηση:
\[ f(x) = x - 2 \]
(Δηλαδή η $C_f$ είναι η ευθεία $y = x - 2$, από την οποία απουσιάζει η τελεία στο $( -2, -4 )$).

Υπολογίζουμε τα καθορισμένα όρια:
- Στο $x = -2$: $\lim_{x\to -2} f(x) = \lim_{x\to -2} (x - 2) = \boldsymbol{-4}$.
- Στο $+\infty$: $\lim_{x\to +\infty} f(x) = \lim_{x\to +\infty} (x - 2) = \boldsymbol{+\infty}$.

\item Διαθέτουμε τη συνάρτηση με σημείο εξόριστο το $x_0 = -2$, ωστόσο το όριο στο σημείο αυτό εκτιμήθηκε ομαλά και συμπεριφέρεται ως απόλυτα πεπερασμένος πραγματικός αριθμός.
Το γεγονός αυτό τη χαρακτηρίζει ως \textbf{εξαιρετέα} (ή \textbf{αιρόμενη}) \textbf{ανωμαλία / σημείο ασυνέχειας}. Μηδενίζεται η αποσύνδεσή της αν η συνάρτηση οριστεί «τεχνητά» στο $-2$ με την τιμή του ορίου της, δημιουργώντας την $f^*(x) = x-2$ για όλο το $\mathbb{R}$.

\item \textbf{Μονοτονία και Ακρότατα στο $\Delta = (-2, +\infty)$:}
Παραγωγίζουμε την $f(x) = x-2$:
\[ f'(x) = (x-2)' = 1 \]
Καθότι $f'(x) = 1 > 0$ για κάθε $x \in (-2, +\infty)$, η συνάρτηση είναι \textbf{γνησίως αύξουσα} σε ολόκληρο το υπό εξέταση διάστημα.
Λόγω της συνεχιζόμενης γνήσιας μονοτονίας, σε ανοικτό διάστημα δίχως κάτοχο μεγίστου (τείνει στο άπειρο) και δίχως τελικό ελάχιστο (ανοιχτό άκρο στο $-2$), η συνάρτηση \textbf{δεν παρουσιάζει κανένα ακρότατο}.

\item Ζητείται ο υπολογισμός του ορισμένου ολοκληρώματος. Επειδή το διάστημα ενσωμάτωσης $[0,2]$ βρίσκεται καθαρά εντός του πεδίου ορισμού:
\begin{align*} 
\dint_0^2 f(x)\,\d x &= \int_0^2 (x - 2)\,\d x \\
&= \left[ \frac{x^2}{2} - 2x \right]_0^2 \\
&= \left( \frac{4}{2} - 4 \right) - (0 - 0) \\
&= 2 - 4 = \boldsymbol{-2}
\end{align*}

\end{erwthma}
\end{thema}

%=========== ΛΥΣΗ - 15ο ΘΕΜΑ Γ ===========
\begin{thema}{Γ}
Δίνεται η εξαρτώμενη $f(x) = x \ln x - x + 1$ στο $x>0$.
\begin{erwthma}

\item \textbf{Μονοτονία και Ακρότατα:}
$f'(x) = (x\ln x)' - 1 + 0 = \ln x + x\frac{1}{x} - 1 = \ln x + 1 - 1 = \ln x$.
- $f'(x) = 0 \Leftrightarrow \ln x = 0 \Leftrightarrow x = 1$.
Η συνάρτηση μειούται (γνησίως φθίνουσα) στο $(0, 1]$ αφού $\ln x < 0$.
Η συνάρτηση αυξάνεται (γνησίως αύξουσα) στο $[1, +\infty)$ αφού $\ln x > 0$.
Διαθέτει \textbf{ολικό ελάχιστο} ακριβώς στο σημείο $\boldsymbol{x = 1}$. Η ποσότητά του καθορίζεται:
\[ f(1) = 1 \cdot 0 - 1 + 1 = 0 \]
Χρησιμοποιώντας τον ορισμό του συνόλου τιμών με ελάχιστο: $\forall x>0 \Rightarrow f(x) \geq f_{\min} \Rightarrow \boldsymbol{f(x) \geq 0}$.

\item \textbf{Εξίσωση Εφαπτομένης στο $x_0 = e$:}
Υπολογισμός σημείου επαφής:
\[ f(e) = e \ln e - e + 1 = e - e + 1 = 1 \]
Κλίση εφαπτομένης:
\[ \lambda = f'(e) = \ln e = 1 \]
Εξίσωση:
\[ y - f(e) = f'(e)(x - e) \Rightarrow y - 1 = 1(x - e) \Rightarrow \boldsymbol{y = x - e + 1} \]

\item Αναζητούμε επίλυση της $x \ln x = x - 1 \Leftrightarrow x \ln x - x + 1 = 0 \Leftrightarrow f(x) = 0$.
Στο ερώτημα 1 αποδείξαμε σαφώς ότι η $f(x)$ είναι παντού μη αρνητική ($f(x) \geq 0$) και ότι η τιμή μηδέν (το απόλυτο ελάχιστο) λαμβάνεται **αποκλειστικά και μόνο** στη θέση $x=1$. Εξαιτίας της γνήσιας μονοτονίας στους εκατέρωθεν κλάδους, η τιμή δεν ανακυκλώνεται ποτέ. 
Οπότε το $x=1$ είναι εξασφαλισμένα η \textbf{μοναδική ρίζα} της εξίσωσης στο $(0, +\infty)$.

\item \textbf{Κινηματική (Ρυθμός Γωνίας):}
Ζητείται να εξεταστεί η εφαπτόμενη γωνία $\omega(t)$ που δημιουργεί η εφαπτομένη της $C_f$ με τον άξονα $x'x$.
Γνωρίζουμε την τριγωνομετρική σύνδεση κλίσης $\lambda$: $\varepsilon\phi(\omega(t)) = f'(x(t))$.
Παραγωγίζοντας ως προς το χρόνο $t$ με τον αλυσιδωτό κανόνα:
\[ \frac{1}{\sigma\upsilon\nu^2(\omega(t))} \cdot \omega'(t) = f''(x(t)) \cdot x'(t) \]
Γνωρίζουμε τη σχέση της τριγωνομετρίας $1 + \varepsilon\phi^2 = 1/\sigma\upsilon\nu^2$. Συνεπώς:
\[ (1 + [f'(x(t))]^2) \cdot \omega'(t) = f''(x(t)) \cdot x'(t) \]

Υπολογίζουμε τα δεδομένα τη χρονική στιγμή επαφής $t_0$ όπου το $x(t_0) = e$:
- Το δεδομένο ταχύτητας $x'(t_0) = 1$.
- Από την πρώτη παράγωγο: $f'(e) = \ln e = 1$.
- Παράγοντας την δεύτερη: $f''(x) = \frac{1}{x}$, οπότε $f''(e) = \frac{1}{e}$.

Αντικαθιστούμε:
\begin{align*}
(1 + 1^2) \cdot \omega'(t_0) &= \frac{1}{e} \cdot 1 \\
2 \cdot \omega'(t_0) &= \frac{1}{e} \\
\omega'(t_0) &= \boldsymbol{\frac{1}{2e} \ \text{rad/sec}}
\end{align*}
Αυτός είναι ο ρυθμός μεταβολής της διαμορφούμενης γωνίας. 
*(Σημείωση: Εάν η άσκηση ζητούσε απλώς γεωμετρικά «τη γωνία» τότε $\varepsilon\phi\omega = f'(e) = 1 \Rightarrow \omega = 45^\circ$, ωστόσο η συμμετοχή του χρόνου $x'(t)$ αποκαλύπτει ερώτημα ρυθμού μεταβολής).*
\end{erwthma}
\end{thema}

%=========== ΛΥΣΗ - 15ο ΘΕΜΑ Δ ===========
\begin{thema}{Δ}
Δίνεται $f:\mathbb{R}\to\mathbb{R}$ με $f(0)=0$ και $f'(x) + f(x) > 0$ για κάθε $x \in \mathbb{R}$.
\begin{erwthma}

\item Έχουμε την βοηθητική συνάρτηση $h(x) = e^x f(x)$. Παραγωγίζοντας:
\[ h'(x) = (e^x)' f(x) + e^x f'(x) = e^x f(x) + e^x f'(x) = e^x \left( f'(x) + f(x) \right) \]
Μας δίνεται ισχυρά το $f'(x) + f(x) > 0$. Επιπλέον η εκθετική είναι πάντα θετική ($e^x > 0$). Το γινόμενο δύο θετικών ποσοτήτων είναι διακριτά θετικό:
\[ \boldsymbol{h'(x) > 0} \]
Οπότε η συνάρτηση $h$ προκύπτει πως είναι \textbf{γνησίως αύξουσα} στο σύνολο του $\mathbb{R}$.

\item Η $h$ είναι γνησίως αύξουσα και παρουσιάζει ένα στρατηγικό σημείο αναφοράς στο $x=0$, όπου:
\[ h(0) = e^0 \cdot f(0) = 1 \cdot 0 = 0 \]
Στηριζόμενοι στη δυναμική της μονοτονίας:
- Για οποιοδήποτε $x > 0 \Rightarrow h(x) > h(0) \Rightarrow e^x f(x) > 0 \Rightarrow \boldsymbol{f(x) > 0}$.
- Για οποιοδήποτε $x < 0 \Rightarrow h(x) < h(0) \Rightarrow e^x f(x) < 0 \Rightarrow \boldsymbol{f(x) < 0}$.

\item \textbf{Μοναδικότητα και αλγεβρική κατανομή μοναδικής ρίζας:}
Βάσει των δεδομένων γνωρίζουμε τον μηδενισμό στην εκκίνηση: $f(0) = 0$. Το $x=0$ αποτελεί \textbf{μία σίγουρη ρίζα}.
Βάσει του Ερωτήματος (2), η συνάρτηση διαθέτει καθαρά αυστηρό θετικό πρόσημο ($f(x)>0$) σε ολόκληρο το δεξί κομμάτι $(0, +\infty)$. Άρα εκεί εξασφαλισμένα δεν περιέχει ρίζες. Το ίδιο συμβαίνει και στο αριστερό κομμάτι $(-\infty, 0)$ όπου είναι σταθερά αρνητική.
Κατά ολοκληρωτική απόδειξη, η εξίσωση $f(x)=0$ φέρει \textbf{μοναδική ρίζα} στο $\mathbb{R}$, αποκλειστικά το $\boldsymbol{x = 0}$.

\item Δίνεται το δεδομένο $f(1) = 1/e$. 
$\bullet$ \textbf{Απόδειξη ότι $\dint_0^1 f(x)\,\d x > 0$:}
Στο διάστημα $(0,1]$, διαπιστώσαμε προηγουμένως (ερώτημα 2) ότι η $f(x) > 0$. Αφού η συνάρτηση συμπεριφέρεται θετικά και δεν ταυτίζεται με τη μηδενική, το ολοκλήρωμά της αποδίδει την απόλυτη επιφάνεια που είναι αυστηρά θετική:
\[ \boldsymbol{\int_0^1 f(x)\,\d x > 0} \]

$\bullet$ \textbf{Εκτίμηση άνω φράγματος:}
Η συνάρτηση $h(x) = e^x f(x)$ είναι γνησίως αύξουσα. Άρα, για κάθε εσωτερικό $x \in [0, 1)$, η $h(x)$ αποκρουστικά παραμένει μικρότερη από το δεξιό άκρο:
\[ h(x) < h(1) \]
Βρίσκουμε το $h(1)$:
\[ h(1) = e^1 \cdot f(1) = e \cdot \frac{1}{e} = 1 \]
Οπότε για τα $x$ στο διάστημα $(0,1)$ ισχύει παντού $e^x f(x) < 1 \Rightarrow f(x) < e^{-x}$.
Εφαρμόζουμε την ολοκλήρωση προς ανίσωση (διετηρώντας το γνήσιο διότι μιλάμε για ανισότητα σε ολόκληρο εσωτερικό πεδίο):
\[ \int_0^1 f(x)\,\d x < \int_0^1 e^{-x}\,\d x = \left[ -e^{-x} \right]_0^1 = -e^{-1} - (-1) = 1 - \frac{1}{e} \]

\begin{parat}{\linewidth}
\textbf{Μαθηματικός Περιορισμός - Τυπογραφικό Εκφώνησης:} Το άνω φράγμα που βρήκαμε είναι $1-\frac{1}{e} \approx 0{,}632$. Η εκφώνηση ζητά να αποδείξουμε την υπερβάλλουσα συνθήκη ότι $\int_0^1 f(x)\,\d x < \frac{1}{e}$ ($0{,}367$), κάτι το οποίο είναι \textbf{άλυτο και μαθηματικά λαθεμένο} γενικότερα. Συνεχείς συναρτήσεις της μορφής αυτής μπορούν κάλλιστα να φτάσουν εμβαδά της τάξης του $0{,}60$ προστατεύοντας επιτυχώς την συνθήκη διαφορικής της υπόθεσης (αφού δεν παρεμποδίζεται το αυξητικό κράμα του $h$). Εξάγουμε το μαθηματικά κομψότερο ορθό όριο. 
\end{parat}

\end{erwthma}
\end{thema}
%=========== ΤΕΛΟΣ ΛΥΣΕΩΝ - ΔΙΑΓΩΝΙΣΜΑ ΠΡΟΣΟΜΟΙΩΣΗΣ 15 ===========
